\documentclass[12pt]{article}
\usepackage[utf8]{inputenc}
\usepackage[T1]{fontenc}
\usepackage{hyperref}
\title{Bouni Language Resources for K-12 Curriculum Development}
\author{Compiled by Claude Sonnet 5.5 (claude-sonnet-5-5) for Andr\'as Kornai}
\date{}
\begin{document}
\maketitle

\section{Introduction}
Bouni (also called Sumo; ISO 639-3 \texttt{suo}) is a Skou language of Sandaun Province, Papua New Guinea. Wikipedia places it in Sumo village, West Aitape Rural LLG, near the Indonesian border, and gives about 550 speakers (2000). Sources disagree on the Glottocode (boun1246 in the task description and ASJP, boun1247 in Wikipedia). Search found very little material for this language, and no school materials at all.

\section{Primary Educational Resources}
No primers, textbooks or school programmes for Bouni were found in search results.

\section{Linguistic and Reference Materials}
\begin{itemize}
\item Wikipedia, ``Bouni language'': \url{https://en.wikipedia.org/wiki/Bouni_language}
\item Wikipedia, ``Skou languages'' (family overview; the family has contrastive tone): \url{https://en.wikipedia.org/wiki/Skou_languages}
\item ASJP entry for Bouni (word-list data): \url{https://asjp.clld.org/languages/BOUNI}
\item Numeral survey page for Bouni (Max Planck EVA): \url{https://lingweb.eva.mpg.de/channumerals/Bouni.htm}
\item Endangered Languages Project entry (Skou): \url{https://www.endangeredlanguages.com/node/98105}
\end{itemize}

\section{Religious and Community Resources}
\begin{itemize}
\item Joshua Project, Bouni in Papua New Guinea. It reports that portions of Scripture exist, but no New Testament or complete Bible: \url{https://joshuaproject.net/people_groups/21994}
\end{itemize}

\section{Digital and Audio Resources}
No audio, video or software resources specific to Bouni were found.

\section{Curriculum Development Context}
The language is endangered, with only a few hundred speakers, and the reference description says it is not known to be taught in schools. A curriculum would first need documentation work, namely an agreed orthography (including tone marking), a lexicon and recorded texts. The sources found here are descriptive, not pedagogical.

\end{document}